% !PRINTED 591
\weberpage{591}
Da die Kongruenz
\[
\frac{2\varrho}{\mu} + \frac{2\varrho_1}{\mu_1} \equiv 1\ (\mathrm{mod}\ 2,\,\omega)
\]
nicht m\"oglich ist, was man ganz wie oben (S.\ 589) zeigt, so ist
f\"ur die Gleichung (3) notwendig und hinreichend:
\weq{4}{\frac{2(\mu_1\varrho - \mu\varrho_1)}{\mu\mu_1} \equiv 0\ (\mathrm{mod}\ 2,\,\omega).}
Daraus folgt:
\weq{5}{\mu\varrho_1 \equiv \mu_1\varrho\ (\mathrm{mod}\ \mu\mu_1),}
und dies ist nur dann m\"oglich, wenn
\weq{6}{\varrho \equiv 0\ (\mathrm{mod}\ \mathfrak{a}), \quad \varrho_1 \equiv 0\ (\mathrm{mod}\ \mathfrak{a}_1).}
Sind umgekehrt die Bedingungen (5) und (6) erf\"ullt, so ist:
\[
\frac{2\varrho}{\mu} - \frac{2\varrho_1}{\mu_1}
\]
eine ganze Zahl und die Gleichung (3) ist befriedigt.

Man nehme nun eine durch $\mathfrak{a}$ teilbare ganze Zahl in $\Omega$:
\weq{7}{\alpha = \mathfrak{a}\mathfrak{c},}
worin $\mathfrak{c}$ relativ prim zu $\mu$ und $\mu_1$ ist. Wenn dann $\varrho$ durch $\mathfrak{a}$
teilbar ist, so kann man eine ganze Zahl $\xi$ nach dem Modul $\mathfrak{m}$
aus der Kongruenz
\[
\varrho \equiv \alpha\xi\ (\mathrm{mod}\ \mu)
\]
bestimmen (Bd.\ II, \S\,166, 7.); dann ist
\weq{8}{\alpha_1 = \frac{\alpha\mu_1}{\mu} = \mathfrak{a}_1\mathfrak{c}}
eine durch $\mathfrak{a}_1\mathfrak{c}$ und durch kein anderes Ideal teilbare ganze
Zahl, und aus (5) ergibt sich
\[
\varrho_1 \equiv \alpha_1\xi\ (\mathrm{mod}\ \mu_1).
\]
5. \gesperrt{Wir erhalten also die gemeinschaftlichen Wurzeln von
$\boldsymbol{A}_\mu(x), \boldsymbol{A}_{\mu_1}(x)$ in der Form:}
\weq{9}{\boldsymbol{s}\!\left(\frac{2\alpha\xi}{\mu}\right) = \boldsymbol{s}\!\left(\frac{2\alpha_1\xi}{\mu_1}\right),}
\gesperrt{worin $\xi$ ein vollst\"andiges Restsystem nach dem Modul $\mathfrak{m}$ durchl\"auft.}

\gesperrt{Hierin kann $\mathfrak{m}$ jedes beliebige ungerade Ideal in $\Omega$ bedeuten.}

Denn man kann zu jedem solchen Ideal zwei Zahlen $\mu, \mu_1$
von der Form (2) w\"ahlen und dann $\omega$ unter den \"aquivalenten
Zahlen so, da\ss{} $\mathfrak{a}$ ungerade und relativ prim zu $\mu\mu_1$ wird.
